% Vietnamese translation for OI Public Library
% Source-PDF: 比赛 CONTEST/2021“MINIEYE杯”中国大学生算法设计超级联赛/2021“MINIEYE杯”中国大学生算法设计超级联赛（10）/2021“MINIEYE杯”中国大学生算法设计超级联赛（10）-题目集.pdf
\documentclass[11pt]{article}
\usepackage{oipl-vn}

\title{Đề bài trận thứ mười HDU Multi-University 2021}
\author{}
\date{}

\begin{document}
\maketitle

\origtitle{Super League of Chinese College Students Algorithm Design 2021, Contest 10}
\sourcepath{contest/hdu-2021-minieye-10-problems.pdf}

\section{Problem A. Pty loves sequence}

\paragraph{Tệp vào:} standard input
\paragraph{Tệp ra:} standard output
\paragraph{Giới hạn thời gian:} 6 seconds
\paragraph{Giới hạn bộ nhớ:} 512 megabytes

Một dãy số nguyên dương \(A\) được gọi là tốt khi và chỉ khi thỏa mãn các
điều kiện sau:
\begin{enumerate}
  \item Độ dài của \(A\) bằng \(n\).
  \item Giả sử \(k\) là giá trị lớn nhất của \(A\). Khi đó \(A\) phải chứa
  \([1,2,3,\ldots,k-1,k]\) như một dãy con.
\end{enumerate}

Dãy con là dãy có thể thu được từ một dãy khác bằng cách xóa một số phần tử
hoặc không xóa phần tử nào, nhưng không thay đổi thứ tự của các phần tử còn
lại.

Trước hết, Pty muốn biết số lượng dãy tốt.

Thứ hai, với mỗi số nguyên \(x\in[1,n]\), Pty cũng muốn biết tổng số lần xuất
hiện của \(x\) trong tất cả các dãy tốt.

Vì các đáp án có thể rất lớn, hãy in các đáp án theo modulo \(p\).

Có tổng cộng \(T\) bộ dữ liệu.

\subsection*{Dữ liệu vào}

Dòng đầu tiên chứa một số nguyên \(T\) (\(1\le T\le 10\)) là số bộ dữ liệu.

Với mỗi bộ dữ liệu, dòng duy nhất chứa hai số nguyên \(n,p\)
(\(1\le n\le 3000,\ 1\le p\le 10^9\)).

\subsection*{Dữ liệu ra}

Với mỗi bộ dữ liệu, in hai dòng.

Dòng đầu tiên chứa một số nguyên: số lượng dãy tốt theo modulo \(p\).

Dòng thứ hai chứa \(n\) số nguyên cách nhau bởi dấu cách. Số nguyên thứ \(x\)
biểu thị tổng số lần xuất hiện của \(x\) trong tất cả các dãy tốt, theo
modulo \(p\).

\subsection*{Ví dụ}

\paragraph{standard input}
\begin{verbatim}
2
3 3
3 1216
\end{verbatim}

\paragraph{standard output}
\begin{verbatim}
0
1 1 1
6
10 7 1
\end{verbatim}

\section{Problem B. Pty with card}

\paragraph{Tệp vào:} standard input
\paragraph{Tệp ra:} standard output
\paragraph{Giới hạn thời gian:} 8 seconds
\paragraph{Giới hạn bộ nhớ:} 512 megabytes

Một trò chơi bài đang trở thành cơn sốt ở Thị trấn Pty. Là thị trưởng, Pty
tất nhiên rất tò mò về trò chơi này và chuẩn bị tổ chức một ván quy mô lớn
trong thị trấn. Trò chơi có luật như sau:
\begin{enumerate}
  \item Có \(N\) người tham gia đứng thành một vòng tròn, mỗi người cầm một
  lá bài.
  \item Chọn ngẫu nhiên một người bắt đầu ở vòng \(1\).
  \item Ở vòng \(i\), giả sử đến lượt Steve và Tom đứng ngay cạnh Steve theo
  chiều kim đồng hồ. Khi đó Tom sẽ lấy bài từ Steve. Nếu \(i\) lẻ, Tom lấy
  \(1\) lá; nếu \(i\) chẵn, Tom lấy \(2\) lá. Nếu lúc này Steve không còn lá
  bài nào, Steve bị xem là bị loại và rời khỏi vòng tròn, còn thứ tự của
  những người khác được giữ nguyên. Ở vòng \(i+1\), đến lượt Tom.
  \item Bất cứ khi nào chỉ còn lại một người tham gia, trò chơi kết thúc.
\end{enumerate}

Hiển nhiên có một số giá trị \(N\) khiến trò chơi trở nên tuần hoàn. Với
những \(N\) như vậy, ta định nghĩa chu kỳ của trò chơi có \(N\) người tham gia
là \(F(N)\). Cụ thể, \(F(N)\) bằng giá trị \(T\) nhỏ nhất sao cho tồn tại một
vòng \(i\), để với mọi \(j\ge i+T\), ở vòng \(j\) và vòng \(j-T\) có cùng tập
người tham gia còn lại và mỗi người có cùng số lá bài. Đặc biệt, nếu trò chơi
có \(N\) người ban đầu sẽ kết thúc, thì \(F(N)=0\).

Thị trấn Pty gồm \(M\) thành phố. Từ mỗi thành phố có thể đi tới bất kỳ thành
phố nào khác thông qua \(M-1\) con đường. Dùng \((x_i,y_i)\) để mô tả con
đường thứ \(i\), nối thành phố \(x_i\) và thành phố \(y_i\). Gọi
\(\operatorname{dist}(x,y)\) là khoảng cách giữa thành phố \(x\) và thành phố
\(y\). Pty nhận thấy rằng nếu chọn tuyến đường từ thành phố \(x\) đến thành
phố \(y\) cho trò chơi, thì số người tham gia là
\[
  p(x,y)=v_x+\operatorname{dist}(x,y).
\]
Với mỗi thành phố \(x\), Pty muốn biết tổng các giá trị chu kỳ khi chọn tất cả
các tuyến đường kết thúc ở thành phố \(x\), tức là bạn cần cho biết
\[
  \sum_{j=1}^{M} F(p(j,x))
\]
với mỗi \(x\).

Cho các giá trị \(M\), \(v_i\), \(x_i\) và \(y_i\), hãy cho Pty biết đáp án.

\subsection*{Dữ liệu vào}

Có tổng cộng \(T\) bộ dữ liệu. \(T\) được cho ở dòng đầu tiên
(\(1\le T\le 10\)).

Với mỗi bộ dữ liệu:

Dòng đầu tiên chứa một số nguyên dương \(M\) (\(1\le M\le 10^5\)) là số thành
phố.

Dòng thứ hai chứa \(M\) số nguyên dương \(v_i\) (\(1\le v_i\le 10^5\)).

Mỗi dòng trong \(M-1\) dòng tiếp theo chứa hai số nguyên dương \(x_i,y_i\),
biểu thị một con đường giữa thành phố \(x_i\) và thành phố \(y_i\)
(\(1\le x_i,y_i\le M\)).

Bảo đảm rằng \(\sum M\le 5\times 10^5\).

\subsection*{Dữ liệu ra}

Có tổng cộng \(T\) dòng.

Với mỗi bộ dữ liệu, in \(M\) số trên một dòng. Số thứ \(i\) là đáp án cho
thành phố \(i\).

\subsection*{Ví dụ}

\paragraph{standard input}
\begin{verbatim}
1
5
1 2 3 4 5
1 2
3 2
3 5
5 4
\end{verbatim}

\paragraph{standard output}
\begin{verbatim}
12 12 0 0 0
\end{verbatim}

\subsection*{Ghi chú}

Với tuyến đường \((2,1)\), số người tham gia là \(3\). Trò chơi diễn ra như
sau:
\[
  (1,1,1)\to(2,1)\to(3),
\]
và trò chơi kết thúc. Do đó \(F(p(2,1))=F(3)=0\).

Với tuyến đường \((5,1)\), số người tham gia là \(8\). Trò chơi diễn ra như
sau:
\[
\begin{aligned}
 &(1,1,1,1,1,1,1,1)\to(2,1,1,1,1,1,1)\to(3,1,1,1,1,1)\\
 &\to(2,2,1,1,1,1)\to\cdots\to(2,2,2,2)\to(4,2,2)\\
 &\to(3,3,2)\to(3,1,4)\to(4,1,3)\to(2,3,3)\to(2,2,4)\\
 &\to(4,2,2)\to(3,3,2)\to\cdots .
\end{aligned}
\]
Vì vậy \(F(p(5,1))=F(8)=6\).

\section{Problem C. Pty loves lines}

\paragraph{Tệp vào:} standard input
\paragraph{Tệp ra:} standard output
\paragraph{Giới hạn thời gian:} 1.5 seconds
\paragraph{Giới hạn bộ nhớ:} 256 megabytes

Bạn cần đặt \(n\) đường thẳng trên một mặt phẳng, bảo đảm không có ba đường
thẳng nào cùng đi qua một điểm và không có hai đường thẳng nào trùng nhau.
Chúng sẽ tạo ra một số giao điểm. Hãy in tất cả các số lượng giao điểm có thể
có.

\subsection*{Dữ liệu vào}

Dòng đầu tiên chứa một số nguyên \(T\) (\(1\le T\le 5\)) là số bộ dữ liệu.

Trong \(T\) dòng tiếp theo, mỗi dòng chứa một số nguyên dương \(n\)
(\(1\le n\le 700\)) là số đường thẳng.

\subsection*{Dữ liệu ra}

In một số dòng; mỗi dòng gồm một số giá trị cách nhau bởi một dấu cách: tất
cả các số lượng giao điểm có thể có.

\subsection*{Ví dụ}

\paragraph{standard input}
\begin{verbatim}
2
3
5
\end{verbatim}

\paragraph{standard output}
\begin{verbatim}
0 2 3
0 4 6 7 8 9 10
\end{verbatim}

\section{Problem D. Pty hates prime numbers}

\paragraph{Tệp vào:} standard input
\paragraph{Tệp ra:} standard output
\paragraph{Giới hạn thời gian:} 3 seconds
\paragraph{Giới hạn bộ nhớ:} 512 megabytes

Pty cho rằng các số nguyên tố là những con số cực kỳ phiền phức. Cụ thể, cậu
ghét \(k\) số nguyên tố đầu tiên.

Theo câu nói ``ghét ai ghét cả đường đi'', nếu một thừa số nguyên tố của số
nguyên \(x\) là số mà Pty ghét, thì Pty cũng ghét \(x\). Bây giờ Pty muốn biết
có bao nhiêu số nguyên trong \([1,n]\) mà cậu không ghét.

Nói hình thức hơn, Pty muốn biết
\[
  \sum_{x=1}^{n} [\,\forall i\in[1,k],\ p_i\nmid x\,],
\]
trong đó \(p_i\) là số nguyên tố thứ \(i\).

\subsection*{Dữ liệu vào}

Dòng đầu tiên chứa một số nguyên \(T\) (\(1\le T\le 10^5\)) là số bộ dữ liệu.

Dòng duy nhất của mỗi bộ dữ liệu chứa hai số nguyên \(n,k\)
(\(1\le n\le 10^{18},\ 1\le k\le 16\)).

\subsection*{Dữ liệu ra}

Với mỗi bộ dữ liệu, in số lượng số nguyên mà Pty không ghét.

\subsection*{Ví dụ}

\paragraph{standard input}
\begin{verbatim}
5
10 1
20 2
30 3
40 4
50 5
\end{verbatim}

\paragraph{standard output}
\begin{verbatim}
5
7
8
9
11
\end{verbatim}

\section{Problem E. Pty loves book}

\paragraph{Tệp vào:} standard input
\paragraph{Tệp ra:} standard output
\paragraph{Giới hạn thời gian:} 10 seconds
\paragraph{Giới hạn bộ nhớ:} 512 megabytes

Pty có một chuỗi \(s\) và một từ điển.

Trong từ điển có \(m\) từ. Mỗi từ \(t_i\) có một giá trị \(w_i\).

Định nghĩa
\[
  f(s,l,r)=\sum_{i=1}^{m} \operatorname{occur}(s,l,r,t_i)\cdot w_i,
\]
trong đó \(\operatorname{occur}(s,l,r,t_i)\) là số lần xuất hiện của chuỗi
\(t_i\) trong \(s[l\ldots r]\).

Giá trị thú vị của chuỗi \(s\) là
\[
  \sum_{l=1}^{|s|}\sum_{r=l}^{|s|} f(s,l,r)^5.
\]

Pty muốn biết giá trị thú vị của chuỗi \(s\).

Vì đáp án có thể rất lớn, hãy in đáp án theo modulo \(10^9+7\).

\subsection*{Dữ liệu vào}

Dòng đầu tiên của dữ liệu vào chứa một số nguyên \(T\) (\(1\le T\le 11\)) là
số bộ dữ liệu.

Với mỗi bộ dữ liệu, dòng đầu tiên chứa một chuỗi \(s\)
(\(1\le |s|\le 5\times 10^5\)).

Dòng thứ hai chứa một số nguyên \(m\) (\(0\le m\le 5\times 10^5\)).

Trong \(m\) dòng tiếp theo, dòng thứ \(i\) chứa một chuỗi \(t_i\) và một số
nguyên \(w_i\)
\[
  \sum_{i=1}^{m}|t_i|\le 5\times 10^5,\qquad 1\le w_i<10^9+7.
\]

Bảo đảm rằng \(s\) và các \(t_i\) đều là chuỗi không rỗng.

\subsection*{Dữ liệu ra}

Với mỗi bộ dữ liệu, in giá trị thú vị của chuỗi \(s\) theo modulo \(10^9+7\)
trên một dòng.

\subsection*{Ví dụ}

\paragraph{standard input}
\begin{verbatim}
1
whkfqbw
3
h 3
hkfq 2
qb 2
\end{verbatim}

\paragraph{standard output}
\begin{verbatim}
75128
\end{verbatim}

\section{Problem F. Pty loves lcm}

\paragraph{Tệp vào:} standard input
\paragraph{Tệp ra:} standard output
\paragraph{Giới hạn thời gian:} 5 seconds
\paragraph{Giới hạn bộ nhớ:} 512 megabytes

\(\operatorname{lcm}(x_1,x_2,\ldots,x_k)\) là bội chung nhỏ nhất của
\(\{x_1,x_2,\ldots,x_k\}\).

\[
  \varphi(x)=\sum_{y=1}^{x}[\gcd(x,y)=1].
\]

Định nghĩa hàm
\[
  f(x,y)=\operatorname{lcm}(x,x+1,\ldots,y-1,y)
  \qquad (x<y),
\]
trong đó \(x,y\) đều là các số nguyên không âm.

Bây giờ Pty muốn biết
\[
  \sum_{x=1}^{+\infty}\sum_{y=x+1}^{+\infty}
  \varphi(f(x,y))\,[\,L\le f(x,y)\le R\,]
  \pmod {2^{32}}.
\]

\subsection*{Dữ liệu vào}

Dòng đầu tiên chứa một số nguyên \(T\) (\(1\le T\le 50\)) là số bộ dữ liệu.

Dòng duy nhất của mỗi bộ dữ liệu chứa hai số nguyên \(L,R\)
(\(1\le L\le R\le 10^{18}\)).

\subsection*{Dữ liệu ra}

Với mỗi bộ dữ liệu, in đáp án theo modulo \(2^{32}\).

\subsection*{Ví dụ}

\paragraph{standard input}
\begin{verbatim}
5
1 10
1 20
20 60
60 100
123456 1234567
\end{verbatim}

\paragraph{standard output}
\begin{verbatim}
5
25
164
160
158548536
\end{verbatim}

\section{Problem G. Pty loves graph}

\paragraph{Tệp vào:} standard input
\paragraph{Tệp ra:} standard output
\paragraph{Giới hạn thời gian:} 10 seconds
\paragraph{Giới hạn bộ nhớ:} 512 megabytes

Cho một đồ thị vô hướng và một chu trình Hamilton của nó, hãy in ra đồ thị đó
có phải là đồ thị phẳng hay không.

Để thuận tiện, các chỉ số đỉnh đã được đánh lại trước sao cho chu trình
Hamilton được cho là
\[
  1-2-3-4-\cdots-(n-1)-n-1.
\]
Các cạnh trên chu trình này không được cho trong dữ liệu vào.

Chú ý rằng đồ thị có thể có khuyên và đa cạnh.

Có tổng cộng \(T\) bộ dữ liệu.

\subsection*{Dữ liệu vào}

Dòng đầu tiên chứa một số nguyên \(T\) là số bộ dữ liệu.

Với mỗi bộ dữ liệu:

Dòng đầu tiên chứa hai số nguyên dương \(N,M\), lần lượt là số đỉnh và số
cạnh ngoài chu trình Hamilton.

Trong \(M\) dòng tiếp theo, mỗi dòng chứa hai số nguyên dương \(x,y\), biểu
thị một cạnh \((x,y)\).

\[
  1\le T\le 26,\qquad
  1\le N,M\le 5\times 10^5,\qquad
  \sum N,\sum M\le 3\times 10^6,\qquad
  1\le x,y\le N.
\]

\subsection*{Dữ liệu ra}

Với mỗi bộ dữ liệu, in \texttt{Yes} hoặc \texttt{No} trên một dòng.

\subsection*{Ví dụ}

\paragraph{standard input}
\begin{verbatim}
2
5 5
1 3
1 4
2 4
2 5
3 5
4 4
1 3
2 4
2 2
1 2
\end{verbatim}

\paragraph{standard output}
\begin{verbatim}
No
Yes
\end{verbatim}

\section{Problem H. Pty loves string}

\paragraph{Tệp vào:} standard input
\paragraph{Tệp ra:} standard output
\paragraph{Giới hạn thời gian:} 10 seconds
\paragraph{Giới hạn bộ nhớ:} 512 megabytes

Pty có một chuỗi \(S\) độ dài \(n\), gồm các chữ cái tiếng Anh thường. Cậu
định nghĩa giá trị của chuỗi \(T\) là số lần xuất hiện của \(T\) trong chuỗi
\(S\).

Bây giờ cậu có \(Q\) truy vấn. Với mỗi truy vấn, cậu cho bạn \(x,y\). Gọi
\(T\) là chuỗi nhận được bằng cách nối tiền tố độ dài \(x\) của chuỗi \(S\)
với hậu tố độ dài \(y\) của nó. Cậu muốn bạn cho biết giá trị của \(T\).

\subsection*{Dữ liệu vào}

Dòng đầu tiên chứa một số nguyên \(T\), là số bộ dữ liệu.

Với mỗi bộ dữ liệu, dòng đầu tiên chứa hai số nguyên \(n,Q\), lần lượt là độ
dài của \(S\) và số truy vấn.

Dòng thứ hai chứa chuỗi \(S\) độ dài \(n\), gồm các chữ cái tiếng Anh thường.

Trong \(Q\) dòng tiếp theo, mỗi dòng chứa hai số nguyên \(x,y\).

\[
  1\le T\le 5,\qquad
  1\le n,Q\le 2\times 10^5,\qquad
  1\le x,y\le n.
\]

\subsection*{Dữ liệu ra}

Với mỗi bộ dữ liệu, in tổng cộng \(Q\) dòng.

Với mỗi truy vấn, in đáp án trên một dòng.

\subsection*{Ví dụ}

\paragraph{standard input}
\begin{verbatim}
1
6 3
ababaa
1 1
2 1
3 1
\end{verbatim}

\paragraph{standard output}
\begin{verbatim}
1
2
1
\end{verbatim}

\section{Problem I. Pty loves SegmentTree}

\paragraph{Tệp vào:} standard input
\paragraph{Tệp ra:} standard output
\paragraph{Giới hạn thời gian:} 15 seconds
\paragraph{Giới hạn bộ nhớ:} 512 megabytes

Pty yêu thích cấu trúc dữ liệu, đặc biệt là cây đoạn.

Pty cho rằng một cây đoạn thỏa mãn tính chất: mỗi nút biểu diễn một đoạn
\([l,r]\). Nút có \(l=r\) được gọi là lá. Một nút khác chọn
\(\mathit{mid}\) trong \([l,r-1]\), rồi lấy \([l,\mathit{mid}]\) và
\([\mathit{mid}+1,r]\) làm hai con.

Pty gán cho mỗi nút trong cây một giá trị. Với lá, giá trị là \(1\). Với nút
có kích thước đoạn của con phải bằng \(k\), giá trị là \(A\). Các nút còn lại
có giá trị \(B\). Pty định nghĩa giá trị của cây là tích các giá trị của mọi
nút.

Pty ký hiệu \(f_n\) là tổng giá trị của tất cả các cây có đoạn của nút gốc là
\([1,n]\). Cậu cho rằng hai cây khác nhau khi và chỉ khi hình dạng của chúng
khác nhau.

Bây giờ Pty có \(Q\) truy vấn. Với mỗi truy vấn, Pty muốn biết giá trị của
\[
  \sum_{i=L}^{R} f_i^2.
\]

Chú ý rằng đáp án rất lớn, bạn chỉ cần tìm đáp án sau khi lấy modulo
\(998244353\).

\subsection*{Dữ liệu vào}

Dòng đầu tiên chứa một số nguyên dương \(T\), là số bộ dữ liệu.

Với mỗi bộ dữ liệu, dòng đầu tiên chứa bốn số nguyên \(Q,k,A,B\).

Trong \(Q\) dòng tiếp theo, mỗi dòng chứa hai số nguyên \(L,R\).

\[
  1\le T\le 5,\qquad
  1\le Q\le 5\times 10^4,\qquad
  0\le A,B<998244353,
\]
\[
  1\le L\le R\le 10^7,\qquad
  1\le k\le 10^7.
\]

\subsection*{Dữ liệu ra}

Với mỗi truy vấn, in đáp án.

\subsection*{Ví dụ}

\paragraph{standard input}
\begin{verbatim}
1
1 1 3 1
4 4
\end{verbatim}

\paragraph{standard output}
\begin{verbatim}
3249
\end{verbatim}

\subsection*{Ghi chú}

Có \(5\) cây khác nhau: \(1\) cây có giá trị \(3\), \(3\) cây có giá trị
\(9\), và \(1\) cây có giá trị \(27\). Tổng là
\[
  3\cdot 1+9\cdot 3+27\cdot 1=57,
\]
và \(57\cdot 57=3249\).

\section{Problem J. Pty plays game}

\paragraph{Tệp vào:} standard input
\paragraph{Tệp ra:} standard output
\paragraph{Giới hạn thời gian:} 5 seconds
\paragraph{Giới hạn bộ nhớ:} 256 megabytes

Pty đang chơi một trò chơi, trong đó cậu chỉ huy một đội lính. Cậu đang đối
mặt với một thử thách trong trò chơi và dự định hoàn thành thử thách sau
\(x\) ngày.

Quái vật cũng chỉ huy một đội lính. Để hoàn thành thử thách, Pty phải đánh
bại những lính này.

Trận chiến diễn ra như sau. Ban đầu, người lính đầu tiên của hai đội sẽ chiến
đấu với nhau. Khi điểm máu của một người lính trở thành \(0\), người lính đó
sẽ được thay bằng người lính tiếp theo trong cùng đội. Một đội được xem là
chiến thắng khi cuối cùng vẫn còn ít nhất một người lính sống, còn đội đối
phương không còn ai sống.

Người lính \(x\) có \(h_x\) điểm máu và có thể gây \(d_x\) sát thương mỗi
giây cho kẻ địch. Chú ý rằng số giây của một trận đấu có thể không phải là
số nguyên.

Cả Pty và quái vật đều liên tục huấn luyện lính của mình, nên điểm máu và sát
thương mỗi giây của mọi người lính sẽ tăng lên mỗi ngày.

Pty muốn biết giá trị nhỏ nhất của \(x\) sao cho cậu hoàn thành được thử
thách nếu nó diễn ra sau \(x\) ngày. Nếu Pty không thể thắng quái vật trong
\(10^{18}\) ngày, hãy in \texttt{none}. Đặc biệt, nếu Pty có thể thắng ngay
bây giờ, đáp án là \(0\).

Chú ý rằng trận chiến sẽ kết thúc trong một ngày, và Pty không thể chọn thứ
tự lính; họ phải chiến đấu theo thứ tự đã cho. Vì vậy, việc duy nhất Pty có
thể làm là chọn ngày tốt nhất để hoàn thành thử thách.

\subsection*{Dữ liệu vào}

Dòng đầu tiên chứa một số nguyên \(T\), là số bộ dữ liệu. Sau đó là \(T\) bộ
dữ liệu.

Trong mỗi bộ dữ liệu, dòng đầu tiên chứa hai số nguyên \(n,m\), lần lượt là
số lính trong đội của Pty và đội của quái vật.

Trong \(n\) dòng tiếp theo, mỗi dòng chứa \(4\) số nguyên. Chúng lần lượt là
điểm máu, mức tăng điểm máu mỗi ngày, sát thương, và mức tăng sát thương mỗi
ngày.

Trong \(m\) dòng tiếp theo, mỗi dòng cũng chứa \(4\) số nguyên. Đây là thông
tin của một người lính trong đội quái vật.

\[
  1\le n,m\le 10^5,\qquad
  \sum(n+m)\le 2\times 10^6,
\]
\[
  0\le h_x,d_x\le 10^6,\qquad
  0\le \text{increment}\le 10^3.
\]

\subsection*{Dữ liệu ra}

Với mỗi bộ dữ liệu, in \texttt{none} hoặc một số nguyên biểu thị giá trị nhỏ
nhất của \(x\) trên một dòng.

\subsection*{Ví dụ}

\paragraph{standard input}
\begin{verbatim}
2
3 2
2 3 1 3
3 2 4 3
1 2 2 3
3 3 4 2
4 1 1 3
1 1
1 1 1 1
1 1 1 1
\end{verbatim}

\paragraph{standard output}
\begin{verbatim}
1
none
\end{verbatim}

\end{document}
